% 页面设置
%%%%%%%%%%%%%%%%%%%%%%%%%%%%%%%%%%%%%%%%%%%%%%%%%%%%%%%%%%%%%%%%%%%%%%
% A4 纸张
\setlength{\paperwidth}{21.0cm}
\setlength{\paperheight}{29.7cm}
% 设置正文尺寸大小
\setlength{\textwidth}{16.1cm}
\setlength{\textheight}{22.2cm}
% 设置正文区在正中间
\newlength \mymargin
\setlength{\mymargin}{(\paperwidth-\textwidth)/2}
\setlength{\oddsidemargin}{(\mymargin)-1in}
\setlength{\evensidemargin}{(\mymargin)-1in}
% 设置正文区偏移量，奇数页向右偏，偶数页向左偏
\newlength \myshift
\setlength{\myshift}{0.35cm}	% 双面打印的奇偶页偏移值，可根据需要修改，建议小于 0.5cm
\addtolength{\oddsidemargin}{\myshift}
\addtolength{\evensidemargin}{-\myshift}
% 页眉页脚相关距离设置
\setlength{\topmargin}{-0.05cm}
\setlength{\headheight}{0.50cm}
\setlength{\headsep}{0.90cm}
\setlength{\footskip}{1.47cm}
% 公式的精调
\allowdisplaybreaks[4]  % 可以让公式在排不下的时候分页排，这可避免页面有大段空白。

% 下面这组命令使浮动对象的缺省值稍微宽松一点，从而防止幅度
% 对象占据过多的文本页面，也可以防止在很大空白的浮动页上放置很小的图形。
\renewcommand{\topfraction}{0.9999999}
\renewcommand{\textfraction}{0.0000001}
\renewcommand{\floatpagefraction}{0.9999}

%%%%%%%%%%%%%%%%%%%%%%%%%%%%%%%%%%%%%%%%%%%%%%%%%%%%%%%%%%%%%%%%%%%%%%
% 字体字号定义
%%%%%%%%%%%%%%%%%%%%%%%%%%%%%%%%%%%%%%%%%%%%%%%%%%%%%%%%%%%%%%%%%%%%%%
% 字号
\newcommand{\yihao}{\fontsize{26pt}{39pt}\selectfont}	  % 一号，1.5  倍行距
\newcommand{\xiaoyi}{\fontsize{24pt}{30pt}\selectfont}  % 小一，1.25 倍行距
\newcommand{\erhao}{\fontsize{22pt}{27.5pt}\selectfont} % 二号，1.25 倍行距
\newcommand{\xiaoer}{\fontsize{18pt}{22.5pt}\selectfont}% 小二，1.25 倍行距
\newcommand{\sanhao}{\fontsize{16pt}{20pt}\selectfont}  % 三号，1.25 倍行距
\newcommand{\xiaosan}{\fontsize{15pt}{19pt}\selectfont} % 小三，1.25 倍行距
\newcommand{\sihao}{\fontsize{14pt}{17.5pt}\selectfont} % 四号，1.25 倍行距
\newcommand{\xiaosi}{\fontsize{12pt}{15pt}\selectfont}  % 小四，1.25 倍行距
\newcommand{\dawu}{\fontsize{10.5pt}{18pt}\selectfont}  % 五号，1.75 倍行距
\newcommand{\zhongwu}{\fontsize{10.5pt}{16pt}\selectfont}% 五号，1.5 倍行距
\newcommand{\wuhao}{\fontsize{10.5pt}{10.5pt}\selectfont}% 五号，单倍行距
\newcommand{\xiaowu}{\fontsize{9pt}{9pt}\selectfont}	   % 小五，单倍行距

\newcommand{\song}{\CJKfamily{song}}
\newcommand{\hei}{\CJKfamily{hei}}
\newcommand{\kai}{\CJKfamily{kai}}
\newcommand{\fs}{\CJKfamily{fs}}
\newcommand{\xkai}{\CJKfamily{xkai}}

% defaultfont 默认字体命令
\def\defaultfont{\renewcommand{\baselinestretch}{1.27}
  \fontsize{12pt}{15pt}\selectfont}

%% 设置目录字体和行间距
\def\defaultmenufont{\renewcommand{\baselinestretch}{1.22}
  \fontsize{12pt}{15pt}\selectfont}

% 固定距离内容填入及下划线
\makeatletter
\newcommand\fixeddistanceleft[2][1cm]{{\hb@xt@ #1{#2\hss}}}
\newcommand\fixeddistancecenter[2][1cm]{{\hb@xt@ #1{\hss#2\hss}}}
\newcommand\fixeddistanceright[2][1cm]{{\hb@xt@ #1{\hss#2}}}
\newcommand\fixedunderlineleft[2][1cm]{\underline{\hb@xt@ #1{#2\hss}}}
\newcommand\fixedunderlinecenter[2][1cm]{\underline{\hb@xt@ #1{\hss#2\hss}}}
\newcommand\fixedunderlineright[2][1cm]{\underline{\hb@xt@ #1{\hss#2}}}
\makeatother

%%%%%%%%%%%%%%%%%%%%%%%%%%%%%%%%%%%%%%%%%%%%%%%%%%%%%%%%%%%%%%%%%%%%%%
% 标题环境相关
%%%%%%%%%%%%%%%%%%%%%%%%%%%%%%%%%%%%%%%%%%%%%%%%%%%%%%%%%%%%%%%%%%%%%%
% 定义、定理等环境
\theoremstyle{plain}
\theoremheaderfont{\hei\bf}
\theorembodyfont{\song\rmfamily}
\newtheorem{definition}{\hei 定义}[chapter]
\newtheorem{example}{\hei 例}[chapter]
\newtheorem{algorithm}{\hei 算法}[chapter]
\newtheorem{theorem}{\hei 定理}[chapter]
\newtheorem{axiom}{\hei 公理}[chapter]
\newtheorem{proposition}[theorem]{\hei 命题}
\newtheorem{property}{\hei 性质}
\newtheorem{lemma}[theorem]{\hei 引理}
\newtheorem{corollary}{\hei 推论}[chapter]
\newtheorem{remark}{\hei 注解}[chapter]
\newenvironment{proof}{\hei{证明}　}{\hfill $\square$ \vskip 4mm}

% 目录标题
\renewcommand\contentsname{\hfill 目　　录 \hfill}
\renewcommand\listfigurename{\hfill 插~图~目~录 \hfill}
\renewcommand\listtablename{\hfill 表~格~目~录 \hfill}
\renewcommand{\bibname}{\hfill 参~考~文~献 \hfill}

%%%%%%%%%%%%%%%%%%%%%%%%%%%%%%%%%%%%%%%%%%%%%%%%%%%%%%%%%%%%%%%%%%%%%%
% 段落章节相关
%%%%%%%%%%%%%%%%%%%%%%%%%%%%%%%%%%%%%%%%%%%%%%%%%%%%%%%%%%%%%%%%%%%%%%
\setcounter{secnumdepth}{4}
\setcounter{tocdepth}{4}
% 设置章、节、小节、小小节的间距
\titleformat{\chapter}[hang]{\normalfont\xiaosan\hei\sf}{\xiaosan\thechapter}{1em}{\xiaosan}
\titlespacing{\chapter}{0pt}{-3ex  plus .1ex minus .2ex}{3.3ex}
\titleformat{\section}[hang]{\sihao\hei\sf}{\sihao\thesection}{1em}{}{}
\titlespacing{\section}{0pt}{0.5em}{0.5em}
\titleformat{\subsection}[hang]{\xiaosi\hei\sf}{\xiaosi\thesubsection}{1em}{}{}
\titlespacing{\subsection}{0pt}{0.5em}{0.3em}
\titleformat{\subsubsection}[hang]{\hei\sf}{\thesubsubsection}{1em}{}{}
\titlespacing{\subsubsection}{0pt}{0.3em}{0pt}
% 缩小目录中各级标题之间的缩进
\dottedcontents{chapter}[0.32cm]{\vspace{0.2em}}{1.0em}{5pt}
\dottedcontents{section}[1.32cm]{}{1.8em}{5pt}
\dottedcontents{subsection}[2.32cm]{}{2.7em}{5pt}
\dottedcontents{subsubsection}[3.32cm]{}{3.4em}{5pt}

% 段落之间的竖直距离
\setlength{\parskip}{1.2pt}
% 段落缩进
\setlength{\parindent}{24pt}
% 定义行距
\renewcommand{\baselinestretch}{1.27}
% 参考文献条目间行间距
\setlength{\bibsep}{2pt}

%%%%%%%%%%%%%%%%%%%%%%%%%%%%%%%%%%%%%%%%%%%%%%%%%%%%%%%%%%%%%%%%%%%%%%
% 页眉页脚设置

\newcommand{\makeheadrule}{%
  \makebox[0pt][l]{\rule[.7\baselineskip]{\headwidth}{0.5pt}}%
  \vskip-.8\baselineskip}

\makeatletter
\renewcommand{\headrule}{%
  {\if@fancyplain\let\headrulewidth\plainheadrulewidth\fi
    \makeheadrule}}

\pagestyle{fancyplain}

\fancyhf{}
\fancyhead[CO]{\song\wuhao{南京信息工程大学硕士学位论文}}
\fancyhead[CE]{\song\wuhao\@ctitle}
\fancyfoot[C,C]{\xiaowu$-$~\thepage~$-$}

% 在上一个空的奇数页清除页眉
\makeatletter
\def\cleardoublepage{\clearpage\if@twoside \ifodd\c@page\else%
  \hbox{}%
  \thispagestyle{empty}%              % 空的页眉
  \newpage%
  \if@twocolumn\hbox{}\newpage\fi\fi\fi}

%%%%%%%%%%%%%%%%%%%%%%%%%%%%%%%%%%%%%%%%%%%%%%%%%%%%%%%%%%%%%%%%%%%%%%
% 列表环境设置
%%%%%%%%%%%%%%%%%%%%%%%%%%%%%%%%%%%%%%%%%%%%%%%%%%%%%%%%%%%%%%%%%%%%%%
\setlist[enumerate]{(1),itemsep=-5pt,topsep=0mm,labelindent=\parindent,leftmargin=*}

%%%%%%%%%%%%%%%%%%%%%%%%%%%%%%%%%%%%%%%%%%%%%%%%%%%%%%%%%%%%%%%%%%%%%%
% 国际单位，以点连接。
%%%%%%%%%%%%%%%%%%%%%%%%%%%%%%%%%%%%%%%%%%%%%%%%%%%%%%%%%%%%%%%%%%%%%%
\sisetup{inter-unit-product = { }\cdot{ }}

%%%%%%%%%%%%%%%%%%%%%%%%%%%%%%%%%%%%%%%%%%%%%%%%%%%%%%%%%%%%%%%%%%%%%%
% 参考文献的处理
%%%%%%%%%%%%%%%%%%%%%%%%%%%%%%%%%%%%%%%%%%%%%%%%%%%%%%%%%%%%%%%%%%%%%%

% \addtolength{\bibsep}{-0.5em}              % 缩小参考文献间的垂直间距
\setlength{\bibhang}{2em}
\bibpunct{[}{]}{,}{s}{}{}

%%%%%%%%%%%%%%%%%%%%%%%%%%%%%%%%%%%%%%%%%%%%%%%%%%%%%%%%%%%%%%%%%%%%%%
%   图形表格
%%%%%%%%%%%%%%%%%%%%%%%%%%%%%%%%%%%%%%%%%%%%%%%%%%%%%%%%%%%%%%%%%%%%%%
\renewcommand{\figurename}{图}
\renewcommand{\tablename}{表}
% \captionstyle{\centering}
% \hangcaption
\captiondelim{\hspace{1em}}
\captionnamefont{\zhongwu}
\captiontitlefont{\zhongwu}
\setlength{\abovecaptionskip}{0pt}
\setlength{\belowcaptionskip}{0pt}


\newcommand{\tablepage}[2]{\begin{minipage}{#1}\vspace{0.5ex} #2 \vspace{0.5ex}\end{minipage}}
\newcommand{\returnpage}[2]{\begin{minipage}{#1}\vspace{0.5ex} #2 \vspace{-1.5ex}\end{minipage}}


%%%%%%%%%%%%%%%%%%%%%%%%%%%%%%%%%%%%%%%%%%%%%%%%%%%%%%%%%%%%%%%%%%%%%%
% 定义题头格言的格式
%%%%%%%%%%%%%%%%%%%%%%%%%%%%%%%%%%%%%%%%%%%%%%%%%%%%%%%%%%%%%%%%%%%%%%

\newsavebox{\AphorismAuthor}
\newenvironment{Aphorism}[1]
{\vspace{0.5cm}\begin{sloppypar} \slshape
    \sbox{\AphorismAuthor}{#1}
    \begin{quote}\small\itshape }
    {\\ \hspace*{\fill}------\hspace{0.2cm} \usebox{\AphorismAuthor}
    \end{quote}
  \end{sloppypar}\vspace{0.5cm}}

% 自定义一个空命令，用于注释掉文本中不需要的部分。
\newcommand{\comment}[1]{}

% This is the flag for longer version
\newcommand{\longer}[2]{#1}

\newcommand{\ds}{\displaystyle}

% define graph scale
\def\gs{1.0}

%%%%%%%%%%%%%%%%%%%%%%%%%%%%%%%%%%%%%%%%%%%%%%%%%%%%%%%%%%%%%%%%%%%%%%%%%%%%%%%%
% 封面摘要
%%%%%%%%%%%%%%%%%%%%%%%%%%%%%%%%%%%%%%%%%%%%%%%%%%%%%%%%%%%%%%%%%%%%%%%%%%%%%%%%
\def\cdegree#1{\def\@cdegree{#1}}\def\@cdegree{}
\def\ctitle#1{\def\@ctitle{#1}}\def\@ctitle{}
\def\caffil#1{\def\@caffil{#1}}\def\@caffil{}
\def\teacher#1{\def\@teacher{#1}}\def\@teacher{} %改这边定义要改5 处，这里3 处，下面1 处，cover里1处
\def\theauthor#1{\def\@theauthor{#1}}\def\@theauthor{}
\def\themajor#1{\def\@themajor{#1}}\def\@themajor{}
\def\direction#1{\def\@direction{#1}}\def\@direction{}
\def\college#1{\def\@college{#1}}\def\@college{}
\long\def\cabstract#1{\long\def\@cabstract{#1}}\long\def\@cabstract{}
\def\ckeywords#1{\def\@ckeywords{#1}}\def\@ckeywords{}
\def\etitle#1{\def\@etitle{#1}}\def\@etitle{}
\long\def\eabstract#1{\long\def\@eabstract{#1}}\long\def\@eabstract{}
\def\ekeywords#1{\def\@ekeywords{#1}}\def\@ekeywords{}

%% 封面
\def\makecover{
  \begin{titlepage}
    \newpage
    \thispagestyle{empty}
    \begin{center}

\begin{table}
\centering
\begin{tabular}{ll}
分类号：\_\_\_\_\_\_\_\_\_\_\_\_\_\_   & ~ ~ ~ ~ ~ ~ ~ ~ ~ ~ ~ ~ ~ ~ ~ ~ ~ 单位代码：\_\_\_\_\_\_\_\_\_\_\_\_\_\_       \\
密~~~~级：\_\_\_\_\_\_\_\_\_\_\_\_\_\_ & ~ ~ ~ ~ ~ ~ ~ ~ ~ ~ ~ ~ ~ ~ ~ ~ ~ 学~ ~ ~ ~ 号：\_\_\_\_\_\_\_\_\_\_\_\_\_\_
\end{tabular}
\end{table}


\begin{figure*}[h]
  \centering
  \includegraphics[width=0.8\textwidth{}]{logo.pdf}
\end{figure*}



      \parbox[t][0.40cm][c]{\textwidth} %这里的0.4控制图片最下边到下面字的最上边的距离
      {
        \begin{center}
          {\xiaoyi\song\textbf{\@cdegree}\\}
          \vspace{1.34cm}
          {\xiaoer\hei\textsf{\@ctitle}\\}
          \vspace{0.13cm}
          {\sanhao\textbf{\@etitle}\\}
        \end{center}
      }

      \parbox[b][8.39cm][c]{\textwidth}
      {
        \vspace{3.8cm} %控制和上一个模块中心行的距离
        \begin{center}
          {
            \xiaosan\song
            \begin{tabular}{p{0.6cm}p{6.4em}@{\extracolsep{0.5em}}lc}
              \vspace{-0.41cm}
              ~ & 申 \hfill 请 \hfill 人 \hfill 姓 \hfill 名：& \fixedunderlineleft[6.08cm]{\@theauthor} & \\
              \vspace{-0.41cm}
              ~ & 指 \hfill 导 \hfill \hfill 教 \hfill 师：& \fixedunderlineleft[6.08cm]{\@teacher} & \\
              \vspace{-0.41cm}
              ~ & 专 \hfill 业 \hfill 名 \hfill 称：& \fixedunderlineleft[6.08cm]{\@themajor} & \\
              \vspace{-0.41cm}
              ~ & 研 \hfill 究 \hfill 方 \hfill 向：& \fixedunderlineleft[6.08cm]{\@direction} & \\
              \vspace{-0.41cm}
              ~ & 所 \hfill 在 \hfill 学 \hfill 院：& \fixedunderlineleft[6.08cm]{\@college} & \\
              \vspace{-0.41cm}
            \end{tabular}
          }
        \end{center}
      }

      \vspace{0.14cm} %这行控制和上一个模块中心的距离

      %{
     	%\xiaoer\xkai
        %南京信息工程大学\\
        %\vspace{0.24cm}
        %\xiaosi\song
       % Nanjing University of Information Science and Technology
      %}
      \xiaoer\song
      \textbf{二〇一八年十一月}

    \end{center}
%    \cleardoublepage %这句话是使偶数页空白，慎用！
  \end{titlepage}
}

%% 编辑独创性说明
\def\originality{
  \renewcommand{\baselinestretch}{1.61}
  \newpage
  \thispagestyle{empty}


%  \begin{center}
    \parbox[t][2.52cm][c]{\textwidth}
    {\xiaoer\song\centerline{\textbf{独创性声明}}}
    \parbox[t][8.68cm][c]{\textwidth}
    {
      \sihao\fs\noindent

~~~~~~本人声明所呈交的论文是我个人在导师指导下进行的研究工作及取得的研究成果。%
本论文除了文中特别加以标注和致谢的内容外，不包含其他人或其他机构己经发表或撰写过的研宄成果，也不包含为获得南京信息工程大学或其他教育机构的学位或证书而使用过的材料。其他同志对本研究所做的贡献均己在论文中作了声明并表示谢意。%

学位论文作者签名：\underline{\hspace{2.4cm}}~~~~~~~ 签字日期：\underline{\hspace{2.4cm}}
}
%  \end{center}

%%编辑使用授权说明
\begin{center}
    \parbox[t][1.52cm][c]{\textwidth}
    {\xiaoer\song\centerline{\textbf{关于论文使用授权的说明}}}
    \parbox[t][8.68cm][c]{\textwidth}
    {
      \sihao\fs\noindent

~~~~~~南京信息工程大学、国家图书馆、中国学术期刊（光盘版）杂志社、中国科学技术信息研宄所的《中国学位论文全文数据库》有权保留本人所送交学位论文的复印件和电子文档，可以采用影印、缩印或其他复制手段保存论文，并通过网络向社会提供信息服务。本人电子文档的内容和纸质论文的内容相一致。除在保密期内的保密论文外，允许论文被查阅和借阅，可以公布（包括刊登）论文的全部或部分内容。论文的公布（包括刊登）授权南京信息工程大学研究生院办理。

学位论文作者签名：\underline{\hspace{2.4cm}}~~~~~~~ 签字日期：\underline{\hspace{2.4cm}}

~~~~~~~~指导教师签名：\underline{\hspace{2.4cm}}~~~~~~~ 签字日期：\underline{\hspace{2.4cm}}
}
  \end{center}

  \renewcommand{\baselinestretch}{1.27}
%  \cleardoublepage %偶数页空白
}





%% 定义摘要的格式
\def\makeabstract{
  \defaultfont
  \chapter*{\hfill 摘　　要 \hfill}
  \addcontentsline{toc}{chapter}{摘　　要} %目录里加入一个章节摘要，但不编号
  \setcounter{page}{1}
  \@cabstract
  \vspace{0.53cm}

  \noindent {\hei{关键词：{\fs\@ckeywords}}}

  \defaultfont
%  \cleardoublepage %偶数页空白
  \chapter*{}
  \addcontentsline{toc}{chapter}{Abstract}
  \vspace{-1.40cm}
  \begin{center}
 %   {\xiaosan\textrm{\@etitle}} %我们南信大英文摘要不需要英文标题

  \end{center}
  \vspace{-0.35cm}
  \begin{center}
    {
      \xiaosan{Abstract}\\
    }
  \end{center}
  \vspace{0.12cm}

  \@eabstract

  \vspace{0.55cm}

  \noindent {\textbf{Key Words:}}~~{\textsf{\@ekeywords}}
 % \cleardoublepage
}


%%
\makeatletter
\def\hlinewd#1{%
  \noalign{\ifnum0=`}\fi\hrule \@height #1 \futurelet
  \reserved@a\@xhline}
\makeatother


% 定义索引生成，产生章节编号？
\def\generateindex
{
  \addcontentsline{toc}{chapter}{\indexname}
  \printindex
 % \cleardoublepage
}




\raggedbottom
